\documentclass[11pt]{article}

% ============================================================================
%  PhiBiz  (working title)
%  Constructive rewrite of the PhiBase arithmetic paper.
%  Body = the core: positions are counts, and exact counting runs backward.
%  Closing chapter = the apparatus, as a study path ("what to learn next").
%  The Athenian (algebra-first) manuscript is pinned separately for later.
%
%  NOTATION (locked): P(n,p) = n + p*phi.  n = unit-step count (integer part),
%  p = golden-step count (phi-coefficient).  Integer part first.
%  Homogeneous form P(n,p,d) = (n + p*phi)/d, d a positive integer.
% ============================================================================

\usepackage{amsmath}
\usepackage{amssymb}
\usepackage{amsthm}
\usepackage[letterpaper,top=1.25in,bottom=1.25in,left=1.25in,right=1.25in,%
            headheight=15pt]{geometry}
\usepackage{booktabs}
\usepackage{array}
\usepackage{enumitem}
\usepackage{fancyhdr}
\usepackage{xcolor}
\usepackage{listings}
\usepackage{parskip}
\usepackage{hyperref}   % loaded last, on purpose

\pagestyle{fancy}
\fancyhf{}
\rhead{\small PhiBase (working title: PhiBiz)}
\lhead{\small Hinds / SpaceStruts Research}
\cfoot{\thepage}

\newtheorem{theorem}{Theorem}[section]
\newtheorem{lemma}[theorem]{Lemma}
\newtheorem{corollary}[theorem]{Corollary}
\theoremstyle{definition}
\newtheorem{definition}[theorem]{Definition}
\theoremstyle{remark}
\newtheorem{remark}[theorem]{Remark}

\newcommand{\Z}{\mathbb{Z}}
\newcommand{\Q}{\mathbb{Q}}
\newcommand{\Zphi}{\mathbb{Z}[\varphi]}
\newcommand{\Qphi}{\mathbb{Q}(\varphi)}
\newcommand{\Pb}[2]{P(#1,\,#2)}
\newcommand{\Pbd}[3]{P(#1,\,#2,\,#3)}
\newcommand{\Norm}{\mathcal{N}}

\lstset{basicstyle=\small\ttfamily, columns=flexible, keepspaces=true}

\begin{document}

\title{%
  \textbf{PhiBase: Exact Arithmetic as a Construction}\\[6pt]
  \large Positions as Counts, and Operations That Run Backward\\[4pt]
  \normalsize\textit{(working title: PhiBiz)}}
\author{%
  James A.~Hinds\\
  SpaceStruts Research\\[4pt]
  \small\texttt{jahbini@icloud.com\quad spacestruts.com}}
\date{July 2026}
\maketitle
\thispagestyle{fancy}

\begin{abstract}
We present an exact integer arithmetic for the golden ratio
$\varphi=(1+\sqrt5)/2$, built from the object up rather than the algebra down.
A position on a line is a record of two step sizes --- a unit step of length
$1$ and a golden step of length $\varphi$ --- and is written as the pair of
counts $\Pb{n}{p}=n+p\varphi$. Addition, subtraction, and multiplication are
closed on these counts (the ring $\Zphi$); division is exact in the field
$\Qphi$, carried by the homogeneous form $\Pbd{n}{p}{d}=(n+p\varphi)/d$, which
keeps its denominator in the open. The organizing property is
\emph{reversibility}: each operation runs backward exactly when its operand is
retained, and the multiplications that reverse with no denominator and without
ever leaving the integer counts are exactly the powers of $\varphi$ --- the
Fibonacci spine. That is the whole of the core. Everything the core supports ---
two lookup tables, a six-count lattice placing dodecahedral geometry in
$\Zphi^3$ with a parity-governed inverse map, a face-plane reflection result
(with the correction that these planes are \emph{not} $H_3$ mirrors), and the
applications to function evaluation, parallel lattice computation, and
exact-arithmetic model inference --- is set out as a study path, with
quantitative claims deferred to measurement and to a companion report.
\end{abstract}

\tableofcontents
\newpage

% ============================================================================
\section{The Construction}
\label{sec:construction}

Ask a narrow question: what can you reach on a line if you are allowed only two
moves --- a step of length $1$ and a step of length $\varphi=(1+\sqrt5)/2$? Every
point you can arrive at is a record of how many of each step you took. That
record is a pair of integers, and the pair \emph{is} the point. Nothing about
this requires a ring to be posited in advance; the algebra is what the counting
turns out to obey, not what it is built on.

The single fact that makes the counts arithmetic rather than mere tally is the
defining identity of $\varphi$,
\[
  \varphi^2 = \varphi + 1,
\]
which reads, in step language, as: two golden steps equal one golden step plus
one unit step. Because of it, any combination of counts --- however multiplied
--- can be rewritten as a plain pair of counts, with no golden step left
squared. Section~\ref{sec:counting} states the representation;
Section~\ref{sec:reverse} states the property the whole paper is organized
around, that these operations run backward exactly when you keep what you moved
by. Everything the core then supports --- the tables, the lattice, the geometry,
and the applications --- is gathered as a study path in
Section~\ref{sec:next}.

The geometry that motivates the choice of $\varphi$ is dodecahedral and
icosahedral: the natural angles are multiples of $36^\circ$, and the vertex
coordinates of these solids, suitably scaled and oriented, are counts of unit
and golden steps~\cite{coxeter,senechal}. The use of $\varphi$ as an arithmetic
base goes back to Bergman~\cite{bergman}; the contribution here is the
constructive reading --- positions as counts, operations as reversible moves ---
and the apparatus that follows from it.

% ============================================================================
\section{Positions Are Counts}
\label{sec:counting}

\begin{definition}[Position]
\label{def:position}
A \emph{position} is a pair of integers $(n,p)$, written $\Pb{n}{p}$ and read as
the point
\[
  \Pb{n}{p} \;=\; n + p\varphi
\]
on the line, where $n$ is the number of unit steps and $p$ the number of golden
steps. Counts may be negative; a negative count is a step taken backward. The
\emph{homogeneous} form $\Pbd{n}{p}{d}=(n+p\varphi)/d$, with $d$ a positive
integer, records a position divided by $d$. The set of positions is
$\Zphi=\{\,n+p\varphi : n,p\in\Z\,\}$~\cite{marcus}; the homogeneous forms
range over $\Qphi$.
\end{definition}

Seven positions recur so often as per-axis coordinate values in dodecahedral
geometry that they earn single-letter names, listed in Table~\ref{tab:decode}.
They are the only values $\{0,\pm1,\pm\varphi,\pm1/\varphi\}$ needed to write the
vertices of the regular solids built from golden triangles, in the coordinate
models used here (Appendix~\ref{app:decode}).

\begin{table}[h]
\centering
\caption{The seven coordinate symbols. (The letter \texttt{P} names the value
  $\varphi$ and is unrelated to the position notation $P(n,p)$; context
  distinguishes them.)}
\label{tab:decode}
\begin{tabular}{ccl}
\toprule
Symbol & $\Pb{n}{p}$ & Decimal value \\
\midrule
\texttt{z} & $\Pb{0}{0}$  & $0$ \\
\texttt{O} & $\Pb{1}{0}$  & $1$ \\
\texttt{o} & $\Pb{-1}{0}$ & $-1$ \\
\texttt{F} & $\Pb{-1}{1}$ & $1/\varphi \approx 0.618$ \\
\texttt{f} & $\Pb{1}{-1}$ & $-1/\varphi \approx -0.618$ \\
\texttt{P} & $\Pb{0}{1}$  & $\varphi \approx 1.618$ \\
\texttt{p} & $\Pb{0}{-1}$ & $-\varphi \approx -1.618$ \\
\bottomrule
\end{tabular}
\end{table}

% ============================================================================
\section{Reversing a Computation}
\label{sec:reverse}

A computation runs forward through the four operations. The question this paper
is built around is when it runs \emph{backward}: given a result and the operands,
when can an input be recovered exactly, with nothing lost on the way? For counts
this is not an existence claim to wave at --- it is a procedure with a cost you
can read off. The reverse of addition is subtraction; the reverse of
multiplication is division; and each reverse is exact \emph{provided you keep the
operand you moved by}. Multiplying $a$ by $b$ and discarding $b$ destroys the
information of which pair produced the result, since many pairs share a product.
Keep $b$, and the road back is open: $a=(a\cdot b)/b$ recovers $a$ with no
rounding, because you divide out exactly what you multiplied in.

Three observations make this concrete. They are the whole of the machinery.

\begin{enumerate}[leftmargin=*]
\item \textbf{Closure you can watch.}
  A count plus a count is a count, and a count times a count is a count.
  Addition is componentwise,
  \[
    \Pb{n_1}{p_1}+\Pb{n_2}{p_2}=\Pb{n_1+n_2}{p_1+p_2},
  \]
  and multiplication, after replacing $\varphi^2$ by $\varphi+1$
  (Appendix~\ref{app:ring}), is
  \[
    \Pb{n_1}{p_1}\cdot\Pb{n_2}{p_2}
      =\Pb{\,n_1 n_2 + p_1 p_2\,}{\,n_1 p_2 + n_2 p_1 + p_1 p_2\,}.
  \]
  Neither result leaves the integer counts. Subtraction, the reverse of
  addition, is componentwise and always exact.

\item \textbf{The costless reversals are the Fibonacci spine.}
  Some multiplications undo without keeping any receipt and without ever leaving
  the counts: exactly those by a position whose reverse is itself a position.
  These are the positions of norm $\pm1$, where
  $\Norm(\Pb{n}{p})=n^2+np-p^2$ (Appendix~\ref{app:ring}), and they are precisely
  the signed powers of $\varphi$,
  \[
    \pm\varphi^k \;=\; \pm\Pb{F_{k-1}}{F_k},\qquad k\in\Z,
    \qquad \Norm(\varphi^k)=(-1)^k,
  \]
  with $F_k$ the $k$-th Fibonacci number ($F_0=0,F_1=1,F_{-1}=1$). This is the
  Fibonacci spine (Section~\ref{sec:next}). Multiplying by $\varphi^k$ is
  reversed by multiplying by $\varphi^{-k}$; on the spine both are one move along
  the index, forward by addition and backward by subtraction. The set of free,
  receipt-less reversals and the central lookup table of the system are the same
  object.

\item \textbf{When the reversal leaves the counts, you keep the divisor.}
  Every other division carries a denominator, and that denominator is the saved
  receipt written in the open. Dividing by $b$ gives
  \[
    \frac{a}{b}=\frac{a\cdot\overline{b}}{\Norm(b)},
    \qquad \overline{\Pb{n}{p}}=\Pb{n+p}{-p},
  \]
  an integer denominator $\Norm(b)$ under an integer-count numerator --- the
  homogeneous form $\Pbd{n}{p}{d}$ of Definition~\ref{def:position}. The
  denominator $d$ is exactly what you multiply back by to undo the division.
  Reducing by $\gcd(n,p,d)$ after each step keeps the receipt in lowest terms.
  Because $\{1,\varphi\}$ is a $\Q$-basis, two triples denote the same field
  element only when they are rationally proportional, so the reduced $d>0$ form
  is a unique representative; $d=1$ says the result is back on the integer
  counts, and $d\ne1$ says how far off, and by what factor, the return trip still
  owes.
\end{enumerate}

\begin{remark}[Why signed counts are needed]
\label{rem:signed}
The non-negative counts $n,p\ge0$ are pure forward motion and grow away from the
origin. Density near zero comes only from the reverse operator: small positions
such as
\[
  |F_k\varphi - F_{k+1}| = \varphi^{-k}
\]
require a step taken backward, e.g. $5\varphi-8=\varphi^{-5}\approx0.09$, which is
$\Pb{-8}{5}$. The signed lattice is counting with subtraction admitted, and it is
what the Dense Grid (Section~\ref{sec:next}) tabulates.
\end{remark}

\subsection{Two Reversals, Worked}

\paragraph{A free reversal (on the spine).}
Take $a=\Pb{1}{1}=1+\varphi=\varphi^2$ and multiply by the spine entry
$b=\Pb{-1}{1}=\varphi^{-1}$. By the product rule,
$a\cdot b=\Pb{0}{1}=\varphi$. To undo it, multiply by $b^{-1}=\varphi=\Pb{0}{1}$:
then $\varphi\cdot\varphi=\Pb{1}{1}=a$, recovered exactly, no denominator, never
off the counts. The receipt was unnecessary because $b$ was a unit.

\paragraph{A reversal that keeps a receipt.}
Divide $a=\Pb{1}{0}=1$ by $b=\Pb{0}{2}=2\varphi$, with $\Norm(b)=-4$ and
$\overline{b}=\Pb{2}{-2}$. The numerator $a\cdot\overline{b}=\Pb{2}{-2}$ over
$-4$ reduces to
\[
  \frac{1}{2\varphi}=\Pbd{-1}{1}{2}=\frac{\varphi-1}{2}\approx0.309,
\]
carrying the receipt $d=2$. To undo it, multiply back by $b=2\varphi$:
$\tfrac{\varphi-1}{2}\cdot2\varphi=(\varphi-1)\varphi=\varphi^2-\varphi=1=\Pb{1}{0}$.
The denominator was the record of what had been divided out; clearing it is the
return trip.

That is the paper's core. The rest is apparatus the core supports.

% ============================================================================
\section{What to Learn Next}
\label{sec:next}

The two tables, the lattice, the geometry, and the applications all sit on the
core of Section~\ref{sec:reverse}. Each is a topic to take up next; below, each
is stated far enough to be usable, with its key result and pointers to the fuller
algebraic development (the companion manuscript) and the measurement work (the
companion report). None of it changes the core.

\subsection{The Fibonacci spine}
\label{ss:spine}

The costless reversals of observation~2, tabulated.

\begin{theorem}[Powers of $\varphi$ are counts]
\label{thm:fibpow}
For all $k\in\Z$, $\varphi^k=F_k\varphi+F_{k-1}=\Pb{F_{k-1}}{F_k}$.
\end{theorem}
\begin{proof}
For $k\ge1$, induction: $\varphi^1=\Pb{0}{1}$, and if
$\varphi^k=\Pb{F_{k-1}}{F_k}$ then $\varphi^{k+1}=F_k\varphi^2+F_{k-1}\varphi
=F_k(\varphi+1)+F_{k-1}\varphi=(F_k+F_{k-1})\varphi+F_k=\Pb{F_k}{F_{k+1}}$.
For $k\le0$, use $F_{-n}=(-1)^{n+1}F_n$, starting from
$\varphi^{-1}=\varphi-1=\Pb{-1}{1}$.
\end{proof}

Each entry is a unit ($\Norm(\varphi^k)=(-1)^k$), so multiplication and division
of exact $\varphi$-powers are index addition and subtraction, and one index step
forward and one back are inverse moves. The spine extends both ways by a single
Fibonacci addition per entry --- the lookup-and-add lineage of logarithmic
number systems~\cite{kingsbury,swartzlander}, but with algebraically exact
entries and no interpolation. The residual scale $\varphi^{-k}$
(Remark~\ref{rem:signed}) reaches $\approx10^{-8}$ near $k=38$, where $F_{38}$ is
on the order of $4\times10^7$; that is a statement about the depth of the spine,
not the resolution of the bounded grid below.

\begin{table}[h]
\centering
\caption{Fibonacci spine entries $\varphi^k=\Pb{F_{k-1}}{F_k}$.}
\label{tab:spine}
\begin{tabular}{rrrl}
\toprule
$k$ & $F_{k-1}$ (unit count) & $F_k$ (golden count) & $\varphi^k$ \\
\midrule
0 & 1 & 0 & $1.000000$ \\
1 & 0 & 1 & $1.618034$ \\
2 & 1 & 1 & $2.618034$ \\
3 & 1 & 2 & $4.236068$ \\
4 & 2 & 3 & $6.854102$ \\
5 & 3 & 5 & $11.090170$ \\
6 & 5 & 8 & $17.944272$ \\
7 & 8 & 13 & $29.034442$ \\
8 & 13 & 21 & $46.978714$ \\
9 & 21 & 34 & $76.013156$ \\
\bottomrule
\end{tabular}
\end{table}

\subsection{The dense grid}
\label{ss:grid}

The spine is sparse and grows away from zero; the dense grid fills the interval
with signed counts, whose small values are the near-cancellations of
Remark~\ref{rem:signed}. For a bound $B$, it is the set of positions $\Pb{n}{p}$
with $|n|,|p|\le B$, sorted by value. Two integer identities are kept per entry
and must not be confused: the \emph{pair index}, the rectangular address
$p\,N_{\text{cols}}+(n-n_{\min})$, and the \emph{value rank}, the position in
sorted order. The counts are exact; only a decimal column is approximate, used
solely for nearest-value binary search.

Sizing and resolution follow from $B$ and are read from the generator, not
assumed. There are $(2B+1)^2$ pairs, distinct in value since $\varphi$ is
irrational; at $B=25$ that is $2601$ pairs, of which $1300$ are positive. The
finest gaps are the residuals $\varphi^{-k}$ reachable with $|F_k|\le B$: at
$B=25$ the smallest is $|8\varphi-13|=\varphi^{-6}\approx0.013$, so the grid
resolves to about $10^{-2}$ there --- not $10^{-8}$, which belongs to the
spine's depth. A function is baked in as a map from input value rank to output
value rank, $L_f[r]$; runtime evaluation is one array read after the input is
encoded.

\subsection{The six-count lattice, and returning by parity}
\label{ss:lattice}

Dodecahedral and icosahedral geometry lifts the counting into three dimensions.
A lattice point is six non-negative counts $[a,b,c,d,e,f]$ along the six
five-fold-axis directions $\mathbf B_1=(\varphi,0,1)$, $\mathbf B_2=(-\varphi,0,1)$,
$\mathbf B_3=(0,1,\varphi)$, $\mathbf B_4=(0,-1,\varphi)$, $\mathbf B_5=(1,\varphi,0)$,
$\mathbf B_6=(-1,\varphi,0)$, giving the forward map into $\Zphi^3$:
\[
  x=\Pb{e-f}{a-b},\qquad y=\Pb{c-d}{e+f},\qquad z=\Pb{a+b}{c+d}.
\]
This is the projection basis of the Kramer--Neri construction of icosahedral
quasicrystals~\cite{kramer}. Writing $x=\Pb{n_x}{p_x}$ and so on, the inverse map
is $a=\tfrac12(n_z+p_x)$, $b=\tfrac12(n_z-p_x)$, $c=\tfrac12(p_z+n_y)$,
$d=\tfrac12(p_z-n_y)$, $e=\tfrac12(p_y+n_x)$, $f=\tfrac12(p_y-n_x)$.

\begin{theorem}[Return by parity]
\label{thm:parity}
A triple $(x,y,z)\in\Zphi^3$ is the forward image of an integer six-count if and
only if $n_z\pm p_x$, $p_z\pm n_y$, and $p_y\pm n_x$ are all even; the six counts
are then uniquely recovered by the inverse map.
\end{theorem}

The forward map takes counts to space and the inverse takes space back to counts;
the $\tfrac12$ in each inverse formula is the receipt, and parity is the test that
it clears --- the three-dimensional face of the reversal principle.
$\varphi$-expansion uses $\varphi\cdot\Pb{n}{p}=\Pb{p}{n+p}$ and returns to an
integer six-count on the even-sum sublattice. The full algebraic treatment is in
the pinned companion manuscript.

\subsection{Reflection in the face planes}
\label{ss:reflect}

Reflect a point through the six face planes of the dodecahedron, whose normals
are the five-fold axes. Each normal has $\mathbf n\cdot\mathbf n=\Pb{2}{1}
=\varphi+2$, with algebraic norm $\Norm(\Pb{2}{1})=5$ and reciprocal
$(3-\varphi)/5$. So reflection maps $\Zphi^3$ into $\tfrac15\Zphi^3$; composed
with the inverse map's factor of two, the round-trip denominator is $10$, and the
reflected point returns to an integer six-count exactly when the numerators (over
$5$) satisfy $N_z\pm P_x,\;P_z\pm N_y,\;P_y\pm N_x\equiv0\pmod{10}$ --- once
again, a receipt-clearing test.

\begin{remark}[These are not mirror planes]
\label{rem:notmirror}
The six face planes are \emph{not} among the fifteen mirror planes of the
icosahedral group $H_3$, whose normals lie along the two-fold (edge-midpoint)
axes; a reflection through a face plane is not a symmetry of the group. The
denominator of $5$ each such reflection introduces is the sign of that --- a true
mirror would permute the vertex set and leave it in $\Zphi^3$. Constructing the
three $H_3$ generators and extending the reflection tables to all fifteen
mirrors, for example to mark hinged edges in a robotic build script, is a
separate construction (Section~\ref{ss:more}).
\end{remark}

\subsection{The lattice as a data-parallel machine}
\label{ss:gpu}

Because neighbour access and reflection are lookups on precomputed address
tables, the geometry maps onto a double-buffered update: each tick dispatches one
pass that reads a node's six neighbours and writes its next state, then swaps
buffers.
\begin{lstlisting}[language=C]
kernel void lattice_step(device Node* last, device Node* next,
                         device int* adj, uint id) {
    next[id] = update(last[adj[id*6+0]], last[adj[id*6+1]],
                      last[adj[id*6+2]], last[adj[id*6+3]],
                      last[adj[id*6+4]], last[adj[id*6+5]]);
}
\end{lstlisting}
The adjacency and reflection buffers are precomputed once; no coordinate
arithmetic runs in the kernel, and the update rule is open (propagation, stress,
routing). Products themselves need no multiplier: a $(2N+1)^2$ table plus adds
returns each product, and at model scale the tables are best shared as a codebook
keyed by distinct quantized value rather than one per weight. Every such table
needs an explicit packed datatype and a proven output range --- the product count
can exceed a signed byte --- and on-chip residency is subject to the device's
threadgroup-memory limit. Throughput is to be measured, not asserted.

\subsection{Inference in exact counts}
\label{ss:ai}

A companion report converts a transformer's stored weights to positions and runs
the matrix arithmetic in exact counts, reproducing reference embeddings on a
tested sample~\cite{phibaseAI}. The scope claim to make here is narrow: the
weight path is exact by construction, because floating-point code cannot run on
tensors stored as integer counts. What the surrounding operations (normalization,
attention, activation) run in is a separate question the companion report
addresses directly.

\subsection{Energy, as an open and measurable question}
\label{ss:energy}

Exact reversible arithmetic retains rather than erases intermediate information.
In the idealized Landauer picture~\cite{landauer,bennett}, information-preserving
operations carry no thermodynamic lower bound on dissipation, unlike the bit
erasure implicit in rounding. Whether that yields measurable savings on real
hardware is empirical: the honest comparison is per-operation energy on the
target device against a floating-point baseline~\cite{horowitz}, and it must be
measured. No fixed multiplier is claimed on paper.

\subsection{Further geometry and layout}
\label{ss:more}

Two open directions. First, the fifteen $H_3$ mirrors: constructing the three
generators (normals perpendicular to the two-fold axes) and extending the
reflection tables to all fifteen planes supports hinged-edge build operations, and
is separate from the face planes of Section~\ref{ss:reflect}. Second, golden-angle
layout: the counts that place polyhedral vertices also place phyllotactic points,
and whether the reversal bookkeeping simplifies golden-angle layouts for arrays
and sensor placement is open.

% ============================================================================
\appendix

\section{Arithmetic on the Counts}
\label{app:ring}
With $\varphi^2=\varphi+1$,
\[
  (n_1+p_1\varphi)(n_2+p_2\varphi)
   =(n_1n_2+p_1p_2)+(n_1p_2+n_2p_1+p_1p_2)\varphi.
\]
The conjugate sends $\varphi\mapsto1-\varphi$, so
$\overline{\Pb{n}{p}}=\Pb{n+p}{-p}$, and the norm
\[
  \Norm(\Pb{n}{p})=\Pb{n}{p}\cdot\overline{\Pb{n}{p}}=n^2+np-p^2\in\Z
\]
is multiplicative. Division uses the conjugate: $a/b=a\,\overline{b}/\Norm(b)$.
A position is a unit ($1/{\Pb{n}{p}}\in\Zphi$) iff $\Norm(\Pb{n}{p})=\pm1$, and
the units are exactly $\pm\varphi^k$.

\section{The Coordinate Vocabulary}
\label{app:decode}
The seven symbols of Table~\ref{tab:decode} span
$\{0,\pm1,\pm\varphi,\pm1/\varphi\}$. In the coordinate models used here, the
dodecahedral and icosahedral vertex sets can be scaled and oriented so that every
coordinate component is one of these seven values; a 3D vertex is then a
three-symbol string (e.g.\ \texttt{OOO}$=(1,1,1)$,
\texttt{zfP}$=(0,-1/\varphi,\varphi)$).

% ============================================================================
\begin{thebibliography}{99}
\bibitem{marcus} D.~A. Marcus, \emph{Number Fields}. Springer, 1977.
\bibitem{bergman} G.~Bergman, A number system with an irrational base,
  \emph{Mathematics Magazine} \textbf{31} (1957), 98--110.
\bibitem{coxeter} H.~S.~M. Coxeter, \emph{Regular Polytopes}, 3rd ed. Dover,
  1973.
\bibitem{kramer} P.~Kramer and R.~Neri, On periodic and non-periodic space
  fillings of $E^m$ obtained by projection, \emph{Acta Crystallographica}
  \textbf{A40} (1984), 580--587.
\bibitem{senechal} M.~Senechal, \emph{Quasicrystals and Geometry}. Cambridge
  University Press, 1995.
\bibitem{kingsbury} N.~G. Kingsbury and P.~J.~W. Rayner, Digital filtering using
  logarithmic arithmetic, \emph{Electronics Letters} \textbf{7} (1971), 56--58.
\bibitem{swartzlander} E.~E. Swartzlander and A.~G. Alexopoulos, The
  sign/logarithm number system, \emph{IEEE Trans.\ Computers} \textbf{C-24}
  (1975), 1238--1242.
\bibitem{landauer} R.~Landauer, Irreversibility and heat generation in the
  computing process, \emph{IBM J.\ Research and Development} \textbf{5} (1961),
  183--191.
\bibitem{bennett} C.~H. Bennett, Logical reversibility of computation,
  \emph{IBM J.\ Research and Development} \textbf{17} (1973), 525--532.
\bibitem{horowitz} M.~Horowitz, Computing's energy problem (and what we can do
  about it), \emph{IEEE ISSCC Digest of Technical Papers} (2014), 10--14.
\bibitem{phibaseAI} J.~A. Hinds, \emph{PhiBase Exact-Count Inference: a
  one-precompute retrofit for transformer models}. SpaceStruts Research, 2026.
  \url{https://spacestruts.com/phibase.html}
\end{thebibliography}

\end{document}
